% !TeX program = lualatex
\documentclass[a4paper,11pt]{article}
\usepackage{luatexja-fontspec}
\usepackage{amsmath,amssymb,amsthm}
\usepackage{mathtools}
\usepackage[unicode]{hyperref}

\newcommand{\PoV}{\operatorname{PoV}}
\newcommand{\dfix}{\bigcirc}
\newcommand{\Rel}{\mathbf{Rel}}
\setlength{\parindent}{0pt}
\setlength{\parskip}{1\baselineskip}

\theoremstyle{definition}
\newtheorem{definition}{定義}[section]
\newtheorem{axiom}[definition]{公理}
\newtheorem{theorem}[definition]{定理}
\newtheorem{remark}[definition]{備考}

\title{冗（Amphigory）の基礎的定式化\\
---視点数・継続残余・無記の関係代数---}
\author{千葉将臣}
\date{2026-07-07}

\begin{document}

\maketitle

\begin{abstract}
本稿では、Freyd と Scedrov による古典的な寓（Allegory）の公理系を基礎とし、そこに視点数（PoV: Point of View）および継続残余 $\rho$ を導入することで、冗の基礎部分を定式化する。寓は、射の交わり、逆関係、およびモジュラー律によって二項関係を抽象化するが、二値截断圏 $\mathcal{B}_{\mathrm{val}}$ に依拠する限り、自己言及的な無限退行や非停止プロセスを未定義として排除する。本稿で導入する冗は、順方向の構成プロセスと逆関係を介した引き戻しプロセスを同時に扱い、証明障害を継続残余として記録する動的代数系である。さらに、非停止プロセスをエラーではなく無記として残余空間へ隔離し、極限において中道不動点 $\dfix$ へ収束させる公理を与える。
\end{abstract}

\section{導入：古典的寓の限界と Amphi- の要請}
P.~J. Freyd と A. Scedrov によって定義された寓は、対象間の二項関係を抽象化した圏であり、射の交わり $\cap$、逆関係 $(-)^\circ$、およびモジュラー律
\[
RS \cap T \subseteq R(S \cap R^\circ T)
\]
を基本構造とする。この枠組みは、集合上の関係の圏 $\Rel$ を圏論的に抽象化するうえで強力である。

しかし、古典的な寓は、射が「成立するか・しないか」という二値截断部分圏 $\mathcal{B}_{\mathrm{val}}$ によって評価される層を基礎としている。この層では、エルブラン空間などにおける動的探索の過程で生じる自己言及的な無限退行、非停止プロセス、あるいは証明探索の発散は、未定義または失敗として扱われる。

冗は、この静的関係の圏に「視点数（PoV）」と「継続残余 $\rho$」の概念を導入し、時間軸に沿った順方向の構成プロセスと、逆関係を介した空間的な引き戻しプロセスの双方向的相互作用を内包する動的代数系である。

また、Amphigory（冗）という語句がもつ「無意味」の含意は、古典論理がエラーとして切り捨ててきたプロセスを、関係代数内で正当に処理するという意味で再解釈される。すなわち、非停止性や評価不能性は単なる失敗ではなく、無記として残余空間へ隔離される。

\section{冗の基本定義と空間構造}

\begin{definition}[評価空間]
冗における評価空間は、少なくとも次の三成分からなる：
\begin{enumerate}
\item 真理値空間 $\mathcal{C}$：四句分別、すなわち是、非、亦是亦非、非是非非の位相を含む空間。
\item 視点数空間 $\mathcal{V} = \mathbb{N} \cup \{\infty\}$：観測に必要な独立した視点の数を表す空間。
\item 残余空間 $\mathcal{L} = \mathbb{R}_{\ge 0}$：評価に伴う物理的・計算的・トポロジカルな摩擦、すなわち証明障害の大きさを記録する空間。
\end{enumerate}
\end{definition}

\begin{definition}[冗の射]
冗 $\mathcal{A}$ において、対象 $X,Y$ 間の射
\[
R: X \to_{\langle p,\rho\rangle} Y
\]
とは、単なる外延的関係ではなく、各ペア $(x,y) \in X \times Y$ に対し、評価推移列を介して次の三つ組を割り当てる多値的な関係である：
\[
(x,y) \longmapsto (c,p,\rho) \in \mathcal{C} \times \mathcal{V} \times \mathcal{L}.
\]
ここで、$c$ は真理値の位相、$p$ は視点数、$\rho$ は継続残余を表す。
\end{definition}

この定義により、古典的な寓における射は、$p$ と $\rho$ が明示的に観測されない二値截断層として回収される。したがって、冗は寓の外側に動的評価の層を付加する拡張である。

\begin{definition}[交わり演算による動的併置]
寓における交わり演算 $R \cap S$ は、外延的な論理積ではなく、観測文脈の動的併置として定義される。

$R$ が $\PoV(R)=n$、$S$ が $\PoV(S)=m$ をもつとき、その交わりは視点数の加算を伴う：
\[
\PoV(R \cap S) = n + m.
\]
同時に、異なる評価文脈の干渉により、継続残余の増大が発生する：
\[
\rho(R \cap S) = \rho(R) + \rho(S) + \Delta_{n,m}.
\]
ここで、$\Delta_{n,m}$ は文脈間の干渉項である。
\end{definition}

交わりをこのように読むことで、「$A$ かつ $B$」という結合は、単なる真理値の積ではなく、複数の観測文脈を同時に保持する操作となる。したがって、論理的結合は計算過程において視点数と残余の変化を伴う。

\begin{definition}[逆関係と逆方向推論]
射 $R$ の逆関係 $R^\circ$ は、ドメインとコドメインを反転させる演算である。このとき視点数は保存される：
\[
\PoV(R^\circ) = \PoV(R).
\]
この演算は、大域的な未来の制約を現在の初期文脈へ空間的に引き戻すための逆方向の推論ベクトルとして機能する。
\end{definition}

Amphi- の語幹が意味する双方向性/両義性は、順方向の構成だけではなく、逆関係を用いた制約の引き戻しを同時に扱う点にある。これにより、評価過程は単なる時系列的展開ではなく、未来制約と現在文脈の衝突を含む双方向的な軌道として記述される。

\section{双方向演繹の代数化：残余モジュラー定理}
古典的なモジュラー律
\[
RS \cap T \subseteq R(S \cap R^\circ T)
\]
は、通常は関係の包含を表す静的な法則として理解される。冗においては、この法則を双方向的な推論軌道の釣り合いと、継続残余 $\rho$ の事前圧縮を保証する原理として読む。

\begin{theorem}[残余モジュラー定理]
冗において、任意の射 $R,S,T$ に対し、以下の操作的包含関係と残余の不等式が成立する：
\[
RS \cap T \subseteq R(S \cap R^\circ T),
\]
\[
\rho_{\mathrm{left}} \ge \rho_{\mathrm{right}}.
\]
ここで、$\rho_{\mathrm{left}}$ は左辺 $RS \cap T$ において発生する大域的な証明障害であり、$\rho_{\mathrm{right}}$ は右辺 $R(S \cap R^\circ T)$ において局所化された証明障害である。
\end{theorem}

\begin{proof}
左辺 $RS \cap T$ は、順方向プロセス $S$ を展開した後に、大域的制約 $T$ との交差検証を行う。このため、制約違反が後段で発見された場合には、事後的な巨大なバックトラックが発生し、その証明障害は $\rho_{\mathrm{left}}$ として蓄積される。

一方、右辺 $R(S \cap R^\circ T)$ は、逆関係 $R^\circ$ を用いて未来の制約 $T$ を空間的に引き戻し、中間プロセス $S$ の展開前に局所制約 $R^\circ T$ として交わりを検証する。これは、順方向推論と逆方向推論の衝突を事前に局所化する操作である。

したがって、右辺では大域的な失敗が発生する前に枝刈りが行われ、残余は局所的に圧縮される。ゆえに、発生する残余について
\[
\rho_{\mathrm{left}} \ge \rho_{\mathrm{right}}
\]
が成り立つ。
\end{proof}

\begin{remark}
この定理は、古典的なモジュラー律を否定するものではない。むしろ、同じ包含式を、計算意味論的には「未来制約の引き戻しによる証明障害の圧縮」として再解釈するものである。
\end{remark}

\section{Nonsense の代数化：無記の残余空間への隔離}
冗という名称の辞書的意味は「Nonsense」あるいは「無意味」である。しかし本稿では、この無意味性を否定的な欠陥としてではなく、古典論理が処理できない評価不能性を代数的に保存するための構造として読む。

\begin{axiom}[無記と極限収束]
冗における自己言及的射の無限列、非停止プロセス、またはフラクタル的な評価軌道において、視点数が発散し
\[
\PoV \to \infty
\]
となるとき、評価は未定義として異常終了するのではなく、無記として残余空間 $\mathcal{L}$ へ隔離される。

この極限操作、すなわち観測者消去により、継続残余 $\rho$ を伴う動的プロセスは、系全体の最大固定点である中道不動点 $\dfix$（空）へと収束する。
\end{axiom}

この公理により、冗は非停止性を単なる失敗として扱わない。むしろ、二値截断圏 $\mathcal{B}_{\mathrm{val}}$ では評価不能となる成分を、無記として残余空間へ移し、体系全体の整合性を保ったまま後続の処理を継続する。

\begin{definition}[中道不動点]
中道不動点 $\dfix$ とは、残余空間を経由した動的プロセスが $E_\infty$ において到達する最大固定点である。単調作用素 $F$ に対して、形式的には
\[
\dfix = \bigvee \{x \mid x \le F(x)\}
\]
として表される。
\end{definition}

ここで $\dfix$ は、単なる停止状態ではない。それは、有限の視点数では決定できないプロセスが、観測者消去を経て到達する極限的な安定状態である。この意味で、$\dfix$ は冗における「空」に対応する。

\section{結語}
本稿では、古典的な寓の公理系に視点数と継続残余を導入することで、冗の基礎的定式化を与えた。

冗の中心的な特徴は、次の三点にまとめられる：
\begin{itemize}
\item 古典的な二値関係を $\mathcal{B}_{\mathrm{val}}$ として含みつつ、その外側に視点数と残余をもつ動的評価層を導入すること。
\item 交わりと逆関係を、観測文脈の併置および未来制約の引き戻しとして再解釈し、双方向演繹を代数化すること。
\item 非停止プロセスや評価不能性を無記として残余空間へ隔離し、極限において中道不動点 $\dfix$ へ収束させること。
\end{itemize}

これにより、冗は、静的な関係代数を、証明探索・非停止性・自己言及・残余の蓄積を含む動的な計算意味論へ拡張する枠組みとして位置づけられる。

\begin{thebibliography}{9}

\bibitem{FreydScedrov1990}
P.~J. Freyd and A. Scedrov.
\newblock \emph{Categories, Allegories}.
\newblock North-Holland, 1990.

\end{thebibliography}

\end{document}
